% Reorganized from the frozen October 15, 2026 supplement.
\section{Shared evidence and rapid/deferred endpoints}
\label{app:reuse}

\subsection{Shared evidence across separated windows}
Let batch $e$ have reference mean $b_e$, $M_e$ candidate means
$b_e-\varepsilon\pm d_e$, and common predetermined precision $\rho_e$.
The final batch contains only its reference; set $M_J=M_{-1}=0$ and
$S_e=\sum_{i\in F_e}q_i(\rho_e)$. Assume
\[
 H_e\ge2S_e,\quad b_J-b_0\le\varepsilon/4,\quad
 d_e+\rho_e+\rho_{e+1}<b_{e+1}-b_e,\quad
 2\max_e\rho_e\le\varepsilon/2.
\]
The algorithm is not told which arm is the reference or when the next window
starts. By induction every batch completes its checkpoints within its first
$2S_e$ rounds, before the next batch appears. On coverage, successive
reference levels satisfy $b_e+\rho_e<b_{e+1}-\rho_{e+1}$. Once batch $e$
completes, the maximum upper endpoint of all available arms is at most
$b_e+\rho_e$. Its candidates are therefore completely classified when
$2\rho_e<d_e$.

An old candidate in batch $e-1$ has upper endpoint at most
$b_{e-1}-\varepsilon+d_{e-1}+\rho_{e-1}$, whereas the new reference has
lower endpoint at least $b_e-\rho_e$. Their difference exceeds
$\varepsilon$, excluding the entire preceding candidate batch without any
further old-arm observations. Older exclusions persist. Every reference has
gap at most $\varepsilon/4$ and total endpoint error at most
$\varepsilon/2$, so all references become positively certified.
If $F=\{e<J:2\rho_e<d_e\}$, the area is consequently at most
\[
 \sum_{e\notin F}M_eH_e+2\sum_{e\in F}M_eS_e+
                2\sum_{e=0}^J(e+1+M_{e-1})S_e.
\]
The terms cover current nonfine candidates, fine candidates, and temporarily
unresolved references plus the preceding candidate batch. All earlier
candidate batches have already been excluded.

For expected regret, every observation with pre-pull count below $q_i$
occurs in its arm's birth window. Such reference observations have zero
regret, and candidate observations cost at most $\varepsilon+d_e$ each.
All observations with pre-pull count at least $q_i$ are ordinary UCB pulls.
For each arm that has a positive gap by $T$, let
$\gamma_i=\min_{t\le T:i\in\mathcal A_t,\,\Delta_i(t)>0}\Delta_i(t)$
and $D_i^T=\max_{t\le T:i\in\mathcal A_t}\Delta_i(t)$.
Use the good-round concentration bound from Section~\ref{app:joint}. At a good UCB round, any positive-gap
pull has pre-pull count at most $12\log(2T)/\gamma_i^2$. Since lifetime counts do not
repeat, this gives the unconditional bound
\begin{equation}
\label{eq:app_reuse_regret}
 ER_T\le\sum_{e<J}\sum_{k=1}^{M_e}(\varepsilon+d_e)q_{e,k}
 +\sum_{i:\gamma_i\text{ exists}}D_i^T
 \left[\left\lceil\frac{12\log(2T)}{\gamma_i^2}\right\rceil-q_i+1\right]_+
 +1.
\end{equation}
The completion-before-next-batch assumption is essential to the first
charging term. It is not an assertion about arbitrary overlapping arrivals.

\subsection{Rapid and deferred endpoints on one family}
\label{app:tradeoff}

Here $J\ge1$, $M\ge\max\{2,J\}$,
$0<\varepsilon\le1/8$, $0<\delta\le1/16$, $Jv\le\varepsilon/4$,
$0<d\le v/32$, and $\eta=d/2$. The reference means are
$b_e=1/2+ev$ and the $M$ candidates born in window $e<J$ have means
$b_e-\varepsilon+\sigma_{e,k}d$. Each window has length $H$, and the final
window contains only a new reference. Set
\[
 K=J(M+1)+1,\quad L_K=\log\frac{4K(K+1)}\delta,
 \quad q_f=\lceil32L_K/d^2\rceil,\quad q_c=\lceil128L_K/v^2\rceil,
\]
\[
 H=128(M+J+1)q_f,\qquad T=(J+1)H,\qquad V=JMH.
\]
These are analysis quantities: the policy uses the selected precision
$\rho_f=d/4$ or $\rho_c=v/8$, not $K,H,T$ or the best-arm identities.
Both choices satisfy the separated-window conditions.

The preceding area bound is at most $B_{J,M}q_f$ for the fine policy and
$V+B_{J,M}q_c$ for the coarse policy, where
\[
 B_{J,M}=2(M+1)\{J(2M+1)+J(J-1)/2\}+2(M+J+1).
\]
For $2\le M$ and $J\le M$, $B_{J,M}\le12JM^2$. Indeed,
$M+1\le3M/2$, $2M+1\le5M/2$, and $J-1\le M$ bound the main term by
$9JM^2$, and the final term by $3JM^2$. The definition of $H$ gives
$B_{J,M}q_f\le V/8$, while $q_c\le q_f$ gives coarse area at most $9V/8$.
Both statements hold on the same fixed-checkpoint coverage event of
probability at least $1-\delta/2$.

For rapid maintenance the area-to-regret lower bound applies. For its upper
bound, every available arm in this family has gap at most $2\varepsilon$.
Hence, on every path, $R_T\le2\varepsilon T$. Since
$(J+1)(M+J+1)/(JM)\le5$,
\[
 2\varepsilon T\le1280\varepsilon JM q_f.
\]
Furthermore $K\le2M^2$, $K+1\le3M^2$ imply
$L_K\le5\log(M/\delta)$, and
$q_f\le165d^{-2}\log(M/\delta)$. Thus
\[
 2^{-20}\frac{\varepsilon JM}{d^2}\log\frac M\delta
 \le\mathfrak R_\Theta(V/8)
 \le2^{18}\frac{\varepsilon JM}{d^2}\log\frac M\delta.
\]
This upper bound uses the stated horizon and must not be extended unchanged
to arbitrarily long windows.

For the coarse policy, within-checkpoint candidate regret is at most
$2\varepsilon JM q_c$. Candidates have
$\gamma_i\ge\varepsilon-d\ge3\varepsilon/4$, $D_i^T\le2\varepsilon$;
old references have $\gamma_i\ge v$, $D_i^T\le Jv$. Dropping the favorable
quota subtraction in \eqref{eq:app_reuse_regret} and absorbing integer terms
gives
\[
 \sup_\Theta ER_T\le2\varepsilon JM q_c+
 64\left(\frac{JM}{\varepsilon}+\frac{J^2}{v}\right)\log(2T)+1.
\]
At fixed $J,M,\varepsilon,v,\delta$, $d\downarrow0$ makes
$T=\Theta(d^{-2})$. Consequently the rapid endpoint is $\Theta(d^{-2})$
and the deferred endpoint is at most $O(\log(1/d))$. No matching lower bound
for the latter endpoint is asserted.
